\chapter*{How to read this working edition}
\addcontentsline{toc}{chapter}{How to read this working edition}

DNA Computing connects algorithms to their molecular realization. A string operation, a mathematical proof and a laboratory measurement can each be useful, but they answer different questions. The book keeps their assumptions and evidentiary roles visible.

This build contains one new chapter. It is designed to be read and challenged,
not to suggest that the remaining planned chapters already exist. Chapter 1
includes a worked example, a bounded proof, a text-sourced vector diagram,
reproducible tables and exercises with concise solution notes.

The source repository separates manuscript, research notes, claims, code and
generated records. Examples are deliberately small enough to inspect. Unit
tests verify stated behaviors; they are not laboratory replication, independent
peer review or evidence of a broad performance advantage.

Read the example first, check its assumptions, then reproduce it. Pay attention
to what would invalidate each conclusion. When an exercise changes an
assumption, predict the result before modifying the program. Unexpected
behavior is a reason to refine the model, not to hide the counterexample.

The earlier edition and the research-reset prospectus remain available in the
repository and Git history. They are not included as chapters in this build.
Final publication still requires independent subject review, a complete
bibliography and index, accessibility review, and confirmed rights metadata.
